\section{Estudo experimental}
\label{sec:estudo-experimental}

O estudo compara quatro configurações para avaliar, de forma progressiva,
o efeito da adaptação do Batch Normalization
(BN)~\cite{ioffe2015batchnorm} e da ativação do Stochastic Depth (SD). A
Tabela~\ref{tab:configuracoes-experimentais} resume as diferenças durante
o treinamento. Em todos os casos, o classificador é treinado e os pesos
das convoluções permanecem congelados. A divisão dos dados, o
pré-processamento e os hiperparâmetros de otimização são mantidos iguais.

\begin{table}[ht]
  \centering
  \caption{Configurações comparadas durante o treinamento.}
  \label{tab:configuracoes-experimentais}
  \begin{tabular}{lccc}
    \hline
    Experimento & \shortstack{Estatísticas\\do BN} &
      \shortstack{Parâmetros\\do BN} & \shortstack{Stochastic\\Depth} \\
    \hline
    Modelo de referência & Fixas & Congelados & Desativado \\
    BN estatístico & Atualizadas & Congelados & Desativado \\
    BN treinável & Atualizadas & Treináveis & Desativado \\
    BN treinável + SD & Atualizadas & Treináveis & Ativado \\
    \hline
  \end{tabular}
\end{table}

O modelo de referência utiliza as estatísticas de normalização do
pré-treinamento. Na configuração BN estatístico, as camadas BN passam a utilizar as
estatísticas dos lotes durante o treino e a atualizar as médias e variâncias
acumuladas, mantendo seus parâmetros de escala e deslocamento congelados.
Na configuração BN treinável, esses parâmetros também são otimizados.
A configuração BN treinável + SD acrescenta o Stochastic
Depth~\cite{huang2016stochasticdepth}, que descarta aleatoriamente a
contribuição de ramos residuais durante o treino. O parâmetro de
probabilidade é definido como $0{,}2$, e a implementação distribui as
probabilidades de descarte de forma crescente com a profundidade dos blocos.

As métricas principais são as acurácias top-1 e top-5. A primeira mede a
proporção de imagens cuja classe correta corresponde à previsão de maior
escore; a segunda considera um acerto quando a classe correta está entre
as cinco previsões de maior escore. As curvas de perda e acurácia de treino
e validação permitem acompanhar a evolução do aprendizado. A avaliação
final também registra precisão e revocação por classe e uma matriz de
confusão, com classes reais nas linhas e preditas nas colunas, para examinar
os erros de classificação.
